\chapter*{Part Synthesis: Foundations}
\addcontentsline{toc}{chapter}{Part Synthesis: Foundations}

A learning task specifies what is observed and what is predicted. A loss describes which prediction errors matter. Their combination defines a population risk,
\[
R(f)=E[\ell(Y,f(X))],
\]
where the expectation belongs to a particular distribution. The training objective replaces this expectation by an observed average. The replacement is useful only when the sampling process and the permitted information match the intended prediction problem.

Linear models express weighted evidence in a chosen representation. Trees, neighbors, and other nonlinear methods change the functions or geometries available to the learner. A more flexible family can reduce approximation error, but the fitted rule must still be selected and evaluated with independent evidence. Dataset design determines which populations, labels, and conditions that evidence represents.

The chapters in this part therefore connect one sequence of questions: what should be predicted, how should error be measured, which rules are plausible, and what observations support a claim about new cases? The next part asks how representations and probability models can be learned rather than fixed in advance.

\clearpage
\begingroup\pagestyle{empty}\cleardoublepage\endgroup
